% THETA3-B18 contribution by /root/definite_bridge, 21 November 2026.
\subsection{THETA3-B18: the remaining definite alternatives have a finite parameter test}\label{B18}

The analytic identity of B17 uses the actual multiplicity plane and
is unconditional on a new global descent. A different question is
whether a prescribed irreducible rank-three constituent also comes
from a totally definite rank-two group. The following statements
make that question more explicit. The global base-change hypotheses
below are inputs, not calculations of the actual A7 constituent.

\paragraph{The complete real deletion table in the canonical case.}
Use B15's real convention, including $\psi(x)=e^{-2\pi i x}$.
Let the rank-three input have Harish--Chandra parameter
$(1,0;-1)$ at $\mu$, and $(1,0,-1;\varnothing)$ at a compact
place. Write the target rank-two splitting character as
$\chi_2(z)=(z/|z|)^{m_0}$, with $m_0$ even; the other character
has any fixed odd exponent $n_0$. A rank-two lift is nonzero
only when $m_0\in\{-2,0,2\}$. Its unique possible signature is
then given by
\[
\begin{array}{c|c|c|c|c}
 m_0&\lambda_\mu-m_0/2&(r_\mu,s_\mu)
       &\lambda_\tau-m_0/2&(r_\tau,s_\tau)\\ \hline
 -2&(2,1;0)&(2,0)&(2,1,0;\varnothing)&(2,0)\\
 0 &(1,0;-1)&(2,0)&(1,0,-1;\varnothing)&(1,1)\\
 2 &(0,-1;-2)&(1,1)&(0,-1,-2;\varnothing)&(0,2)
\end{array}
\tag{B18a}
\]
All six listed lifts are nonzero. In particular total definiteness
requires $m_{0,\mu}\in\{-2,0\}$ and
$m_{0,\tau}\in\{-2,2\}$ elsewhere. The other splitting character
changes the output by a determinant twist and does not change this
table.

\textbf{Proof.}
Ichino's Theorem 4.1(ii), pp.11--12, with $n=3,m=2$, requires a
zero in the shifted parameter: its condition $p_1+q_1>0$ is
necessary. This proves the asserted restriction on $m_0$.
Delete that zero. Each remaining positive entry in the first
block and negative entry in the second block contributes to the
positive target block; the other entries contribute to the
negative target block. The same theorem gives exactly the six
signatures in (B18a), so no other rank-two signature is possible.
For nonvanishing, apply Atobe's Definition 1.6 and Theorem 1.7(2),
pp.4--5. In five cases the zero chain has length one and the
listed signature is its base point $(r_\lambda,s_\lambda)$.
For the compact input with $m_0=0$, the maximal chain is instead
$(1,0,-1)$ of length three; its base point is $(0,0)$, and the
listed $(1,1)$ is obtained by the theorem's equal increment
$l=1$. This checks that case without mistaking its chain length
for one. The output Harish--Chandra parameter is obtained from
the retained entries by adding $n_0/2$, as in B15.
\hfill$\square$

The conversion between directions is explicit. Start with a
Hermitian rank-three $V$ and skew-Hermitian $W$ of rank two.
For a trace-zero scalar $\delta$, switch to the Hermitian space
$(W,\delta^{-1}h_W)$ and the skew-Hermitian space
$(V,\delta h_V)$. The tensor symplectic form is unchanged.
Ichino, \S2.2, p.5, states that this switches the two factors of
the same oscillator. At a real place take $\delta=i$, or a
positive multiple of it in the fixed complex embedding, so the
signatures in (B18a) are precisely those of the Hermitian
multiplicity space. The splitting characters switch their names;
none is inverted in this factor-switch rule. These identifications
must be made before applying the table to the archive convention.

\paragraph{The exact finite-place parameter and inner-form rule.}
Now label the original pair by $V_3$ Hermitian and $W_2$
skew-Hermitian, with characters $\chi_3|_{\A_F^\times}=\epsilon$
and $\chi_2|_{\A_F^\times}=1$. In the local quotient convention
\[
 \Hom_{\U(V_3)\times\U(W_2)}
       (\omega,\Pi\boxtimes\Sigma)\ne0,
\]
Gan--Ichino, \emph{The Gross--Prasad conjecture and local theta
correspondence}, arXiv:1409.6824v2, Theorem 4.4, pp.14--15,
gives at every nonsplit finite place
\[
 \phi_{\Pi_v}
    =(\phi_{\Sigma_v}\otimes\chi_{3,v}^{-1}\chi_{2,v})
                                      \oplus\chi_{2,v}.
 \tag{B18b}
\]
This theorem covers all local $L$-parameters, not just tempered
ones. For fixed $\Pi_v$ and fixed splitting characters, the
smaller parameter is obtained by deleting one occurrence of
$\chi_{2,v}$. The theorem's packet bijections then determine a
unique smaller representation, including its Hermitian inner
form. If the smaller parameter does not contain $\chi_{3,v}$,
use Theorem 4.4(i)(b) for the fixed target sign; if it does,
use Theorem 4.4(ii)(c). The enhanced packet character is the
restriction prescribed in \S4.6, proved by Theorem 1.3. This
last assertion is a finite-place statement; (B18a) supplies the
real-place assertion separately.

The same parameter formula holds at split finite places, where
there is a unique space of each dimension. One can check the
convention rather than silently omit a contragredient.
M\'inguez, \emph{Correspondance de Howe explicite: paires duales
de type II}, Ann. Sci. ENS 41 (2008), Corollary 6.3, p.734,
gives $\phi^\vee\oplus1$ for his untwisted matrix model in
ranks $2\to3$. His equation (1.4) and footnote on p.721 state
that replacing the first action by transpose inverse replaces
the input by its contragredient. Thus the switched model has
parameter $\phi\oplus1$. Twisting its first action by
$\chi_3\circ\det$ and its second by $\chi_2\circ\det$ gives
(B18b). This is the split normalization in Ichino's \S8.1.3,
Lemma 8.4, p.36; M\'inguez's statement also covers non-tempered
inputs. In particular (B18b) is available at every unramified
place, without a Ramanujan assumption.

\paragraph{A conditional global uniqueness test.}
Suppose that a fixed irreducible automorphic $\Pi$ on $\U(V_3)$
has the real components specified above and occurs in a global
theta pair with an irreducible automorphic $\Sigma$ on $\U(W_2)$.
Here $\Sigma$ labels the local oscillator quotient in the
preceding display. In the archive bilinear integral of B17, the
integrating automorphic input has type $\Sigma^\vee$, or its
chosen unitary conjugate realization. Assume that the standard
weak base changes $\mathcal B_3$ and $\mathcal B_2$ are specified isobaric
sums of unitary cuspidal automorphic representations of total
ranks three and two over $E$, with the local parameters in (B18b)
at almost all places. Then
\[
 \mathcal B_3=
   (\mathcal B_2\otimes\chi_3^{-1}\chi_2)\boxplus\chi_2.
 \tag{B18c}
\]
Indeed (B18b) identifies the unramified parameters at almost all
places, and strong multiplicity one for unitary isobaric
representations gives (B18c). The version being used is the
usual finite-exception case of Rajan, \emph{On strong
multiplicity one for automorphic representations},
arXiv:math/0210235, Theorem 1, p.2, and its isobaric extension
stated after the proof on p.6.

If $\mathcal B_2$ is cuspidal on $\mathrm{GL}_2$, then
$\chi_2$ is the unique one-dimensional constituent of the
isobaric cuspidal support in (B18c). Consequently any other
rank-two global realization of this same $\Pi$, in the same
standard-parameter convention and with base changes of the
same specified isobaric kind, must have the same rank-two character
$\chi'_2=\chi_2$. To prove this, apply (B18c) to the second
realization and compare the unique isobaric cuspidal supports.
Its displayed $\mathrm{GL}_1$ summand must be the unique one
in the first support.

Changing $\chi_3$ while fixing $\chi_2$ only twists the smaller
representation by the corresponding determinant character.
The oscillator action on the smaller group is multiplied by
that character, whereas its action on $\U(V_3)$ is unchanged.
Thus the local space is unchanged. More explicitly, write
$\chi'_3/\chi_3=\eta$, and use Hilbert 90 to define
$\nu_\eta(z/\bar z)=\eta(z)$. The source representation is
transported by $\nu_\eta\circ\det$; its parameter in (B18b)
is multiplied by $\eta$, exactly cancelling $\chi_3'^{-1}$.
The finite-place packet bijection and the real table then
force the original smaller local isometry classes.
Therefore, under these base-change and cuspidality hypotheses,
a mixed real place of the original smaller space obstructs a
totally definite rank-two realization of the same $\Pi$.
This is a conditional obstruction for the specified constituent;
no cuspidal $\mathrm{GL}_2$ transfer has been computed here for
an actual A7 wedge.

A compensating determinant twist of the rank-three realization
does not remove this obstruction. It multiplies every summand
of $\mathcal B_3$ and the distinguished $\chi_2$ by the same
character. At a real place it adds $a$ to every entry of
$\lambda$ and adds $2a$ to $m_0$, leaving
$\lambda-m_0/2$ and hence (B18a) unchanged. This recovers the
invariance proved in B16 at the level of the explicit table.

If instead $\mathcal B_2=\eta_1\boxplus\eta_2$ is a sum of
characters, (B18c) has at most three distinct character
constituents. Any alternate $\chi'_2$ must be one of them,
must have the allowed restriction to $\A_F^\times$, and must
pass (B18a) at every real place. For each candidate, deletion
in (B18b) and its enhanced packet rule determine the finite
local smaller representation. Coherence and global occurrence
remain to be proved for that candidate, for example through
the pole criterion of B16. This is a finite list of specified
character tests, not a claim that every deletion occurs globally.
Nontrivial Arthur $\mathrm{SL}_2$ blocks or base changes outside
the unitary isobaric hypothesis require their own analysis;
they are not covered by the preceding two cases.

\paragraph{An equal-rank central-value implication.}
There is also a precise sufficient condition for constructing a
new definite rank-two group from rank-two input. Use exactly
Jiang--Wu's notation in \S3.1--3.2 and Theorem 3.3(1), p.17.
Let $X$ be skew-Hermitian of rank two, let $\sigma$ be irreducible
cuspidal on $\U(X)$, and let $\chi$ be unitary with
$\chi|_{\A_F^\times}=1$.
Put $\rho=\sigma\otimes\chi^{-1}$ in their unitary-group
character convention. Suppose
\[
 L^S(1/2,\sigma\times\chi)\ne0,
 \tag{B18d}
\]
with a finite value, suppose $\rho$ has at least one local
component which is not one dimensional, and suppose that
at each real place every nonzero rank-two local theta lift
of the indicated input is definite. Then some global definite
Hermitian $Y$ of rank two has a nonzero cuspidal theta lift of
$\rho$. Its finite local forms are coherent automatically.

\textbf{Proof.}
In that theorem take $m=2$, $\epsilon_\chi=0$ and $j=1$.
Then $m-\epsilon_\chi$ is even and
$(m+1-\epsilon_\chi)/2-j=1/2$. Hence
$\operatorname{LO}_{\psi,\chi}(\rho)\leq2$.
The dimension is zero or two. The rank-zero theta series is
a determinant character, as in Jiang--Wu's \S3.1; the local
hypothesis excludes its occurrence. The first occurrence is
therefore two and is nonzero cuspidal by the definition of
lowest occurrence. The assumed real local condition makes its
global space definite. The global construction supplies
coherence at the finite places. \hfill$\square$

The real condition in this last statement can itself be made
explicit when the fixed input $\rho$ is discrete series at every
real place. For a mixed input $\U(1,1)$ with parameter $(a;b)$,
$a,b\in\Z+1/2$, $a\ne b$, choose the even exponent $m_0$ of
the target splitting character so that $m_0/2$ lies strictly
between $a$ and $b$. The equal-rank case $k=0$ of Ichino's
Theorem 4.1(ii) then gives the unique signature $(2,0)$ if
$a>b$, and $(0,2)$ if $a<b$. Atobe's base-point criterion
proves nonvanishing. For a compact input, choose $m_0/2$
outside the interval spanned by its two half-integral parameters;
the same rules again give a unique definite target and a
nonzero lift. Suitable integers exist in all these cases.
B15's character-extension lemma globalizes these even infinity
exponents, allowing finite ramification. If the desired input
is fixed as $\rho$, set $\sigma=\rho\otimes\chi$ so that
$\rho=\sigma\otimes\chi^{-1}$ holds in Jiang--Wu's convention.
Condition (B18d) must then be tested for this resulting pair
$(\sigma,\chi)$; the freedom to choose infinity characters
does not prove the central value is nonzero.

At a nonsplit finite place the equal-rank local test is also
explicit. In Gan--Ichino's conventions, their Theorem 4.1(ii),
p.13, says
\[
 \begin{gathered}
 \theta_{Y_v}(\rho_v)\ne0
 \quad\Longleftrightarrow\quad
 \epsilon\bigl(1/2,\phi_{\rho_v}\chi_{Y,v}^{-1},
                      \psi_{2,v}^E\bigr)
               =\epsilon(Y_v)\epsilon(X_v),\\
 \psi_{2,v}^E(x)=\psi_v(\operatorname{Tr}_{E_v/F_v}(\delta x)).
 \end{gathered}
 \tag{B18e}
\]
Its two sides determine one of the two finite inner forms.
The additive character in this epsilon factor has no factor
$1/2$; the $\psi^E$ used for the Whittaker normalization in
that paper does have that factor. These characters must not be
confused. When (B18d) applies, the global lowest-occurrence
theorem constructs a coherent collection of these finite forms;
otherwise guessing their signs separately does not prove
coherence or a nonzero global lift.

This central-value condition is different from B15's upward
value-at-one test and B16's rank-three pole-at-one test. It can
produce an equal-rank definite transfer when its hypotheses are
verified. It does not say that the resulting smaller
representation lifts back to the prescribed $\Pi$, or that it
satisfies the arithmetic kernel/covolume equation of B15.
In particular the conditional uniqueness obstruction above
must still be respected when trying to use such a transfer in
a rank-two-to-rank-three realization of the original wedge.
